\section{Implement a Queue using Stacks}
\label{sec:sq-queue-using-stack}

\problemstatement{Implement a queue (FIFO: \textit{enqueue},
\textit{dequeue}, \textit{front}) using only the standard operations of a
stack (\textit{push}, \textit{pop}, \textit{top}) as building blocks.}

\begin{patternbox}
The reverse simulation of Section~\ref{sec:sq-stack-using-queue}. The key
fact to recall: reversing a sequence \textbf{twice} restores its original
order -- so pouring elements through \emph{two} stacks in succession
(stack $\to$ stack $\to$ stack) flips LIFO order into FIFO order, because
each individual pour through a stack reverses the order once.
\end{patternbox}

\subsection*{Understanding the problem carefully}

A single stack inherently reverses whatever order you feed into it: if
you push $1,2,3$ and then pop everything, you get $3,2,1$ -- reversed. If
we use \textbf{two} stacks, call them \textit{inStack} and
\textit{outStack}, and whenever \textit{outStack} runs empty we pour all
of \textit{inStack} into it (popping from \textit{inStack} and pushing
each onto \textit{outStack}), the act of pouring reverses the order a
\emph{second} time -- and two reversals cancel out, restoring the
original (FIFO) order on \textit{outStack}'s pop sequence.

\begin{ideabox}[Two Stacks: One for Input, One for Output]
$\textit{enqueue}(x)$: simply push $x$ onto \textit{inStack} -- $O(1)$,
no reordering needed yet. $\textit{dequeue}()$ (and $\textit{front}()$):
if \textit{outStack} is empty, pour the entirety of \textit{inStack} into
it (pop each element off \textit{inStack}, push it onto \textit{outStack});
this reverses \textit{inStack}'s order, placing the
\emph{oldest-enqueued, not-yet-dequeued} element on top of
\textit{outStack}. Then pop (or peek) \textit{outStack}'s top.
\end{ideabox}

\subsection*{Why pouring only when \textit{outStack} is empty gives amortized $O(1)$}

If we poured on \emph{every} dequeue, each dequeue would cost $O(n)$.
Instead, we only pour when \textit{outStack} is completely empty -- and
once poured, \textit{outStack} correctly holds a whole batch of elements
in the right FIFO order, which can each be dequeued in $O(1)$ without
needing another pour until that batch is exhausted. Each individual
element is poured (moved from \textit{inStack} to \textit{outStack}) at
most \emph{once} in its entire lifetime, so across any sequence of $n$
enqueue/dequeue operations, the total cost of all pours combined is
$O(n)$, giving $O(1)$ \emph{amortized} time per operation even though some
individual dequeues are more expensive than others.

\subsection*{Algorithm}

\begin{enumerate}
  \item $\textit{enqueue}(x)$: push $x$ onto \textit{inStack}.
  \item $\textit{dequeue}()$ / $\textit{front}()$: if \textit{outStack}
        is empty, repeatedly pop from \textit{inStack} and push onto
        \textit{outStack} until \textit{inStack} is empty. Then pop (or
        peek) \textit{outStack}.
\end{enumerate}

\subsection*{Dry run}

\dryrun{Enqueue $1,2,3$; dequeue twice; enqueue $4$; dequeue once.}

\begin{itemize}
  \item Enqueue $1,2,3$: \textit{inStack} (bottom to top) $=[1,2,3]$;
        \textit{outStack} $=[]$.
  \item Dequeue: \textit{outStack} empty, so pour: pop $3,2,1$ from
        \textit{inStack}, push onto \textit{outStack}, giving
        \textit{outStack} (bottom to top) $=[3,2,1]$. Pop \textit{outStack}'s
        top: returns $1$. \textit{outStack}$=[3,2]$.
  \item Dequeue: \textit{outStack} not empty ($[3,2]$). Pop top: returns
        $2$. \textit{outStack}$=[3]$.
  \item Enqueue $4$: push onto \textit{inStack}$=[4]$.
        \textit{outStack} still $=[3]$, untouched.
  \item Dequeue: \textit{outStack} not empty. Pop top: returns $3$.
\end{itemize}

The dequeue sequence $1,2,3$ correctly matches the enqueue order.

\subsection*{C++ implementation}

\begin{lstlisting}
#include <bits/stdc++.h>
using namespace std;

class QueueUsingStack {
    stack<int> inStack, outStack;

    void transferIfNeeded() {
        if (outStack.empty()) {
            while (!inStack.empty()) {
                outStack.push(inStack.top());
                inStack.pop();
            }
        }
    }

public:
    void enqueue(int x) {
        inStack.push(x);
    }

    int dequeue() {
        transferIfNeeded();
        int val = outStack.top();
        outStack.pop();
        return val;
    }

    int front() {
        transferIfNeeded();
        return outStack.top();
    }
};

int main() {
    QueueUsingStack q;
    q.enqueue(1); q.enqueue(2); q.enqueue(3);
    cout << q.dequeue() << " " << q.dequeue() << endl;
    q.enqueue(4);
    cout << q.dequeue() << endl;
    return 0;
}
\end{lstlisting}

\begin{complexitybox}
\textbf{Time Complexity.} $\textit{enqueue}$ is always $O(1)$.
$\textit{dequeue}$ and $\textit{front}$ are $O(1)$ \textbf{amortized}
(occasionally $O(n)$ for a pour, but each element is poured at most once
across its lifetime).
\textbf{Space Complexity.} $O(n)$ across the two stacks combined.
\end{complexitybox}

\begin{takeawaybox}
Reversing twice restores original order -- this is the entire idea behind
using two stacks to simulate a queue, and it is worth remembering as a
general trick: any time an algorithm needs to ``undo'' a reversal it was
forced into, pouring through a second stack is often the simplest fix.
\end{takeawaybox}
